% The scoped proofs are pinned by the November 9 unit-pair research receipts.
\section{Joint moments for the wide cutoff pairing}
\label{pca:joint-core}

We work at the pure scales \(A=A_1\asymp R\asymp D^{3/4}\),
\(A_2\asymp1\), with \(U^2\le D\le 7U^2/2\). The rows are
the original unweighted nonprincipal rows, with fixed \(S,\nu,u_0\),
their literal ray and capacity restrictions, and their primitive
conductors \(Q_u\asymp R\). Use one smooth complex \(W\) supported
on \([1,2]\), the full ordered cutoff, and the original raw zeros.
A target bound means
\(O((DH)^\epsilon\mathcal M_{24}(W)^2 D^{3/4})\) for the
specified contribution to the normalized pairing \eqref{pca:cross}.
Constants may depend on the fixed arithmetic data, scale constants
and \(\epsilon>0\). Put \(E=Nd\) and \(e=\log_D E\).
Here \(d\) is the intermediate divisor from the exact Möbius
identity. All estimates include the original middle equality subtraction.

\begin{theorem}[Full-frequency partner unit part]\label{pca:joint-full}
Take the unit part of \(C_d\) on the partner side of
\eqref{pca:wide-factor}, with all partner frequencies including the
re-expressed direct term. A selection of the actual ordered pairs
\(mn=d\) may depend on \(d,m,n\). Write \(L_1=U/Nm\),
\(L_2=U/Nn\) and \(l_s=\log_D\max(1,\min(L_1,L_2))\).
The target holds over all \(E\lesssim D^{1/6}\), and more generally
on the dyads with \(e\le1/4\) and \(l_s\ge1/4+e/2\) at
the level of powers. The remaining upper cost at a single
dyad with \(e\le1/4\), relative to \(R/E\), has exponent
\((1/4+e/2-l_s)_+/6\), where \(x_+=\max(0,x)\).
The all-divisor bound for this component has power \(37/48\).
\end{theorem}

The low-divisor estimate uses the whole wide product support, even
for balanced pairs with bounded nonunit \(d\). Its gain comes from
the forced radical of the actual product ideal in the joint positive
norm. For instance, at \(e=1/4\) with the longer length of order
\(D^{1/2}\), the upper comparison still costs \(D^{1/48}\)
above \(R/E\). This excess describes the estimate.

\begin{theorem}[Selected partner frequencies]\label{pca:joint-modes}
Let \(T\le N\mathfrak k<2T\) be an actual partner frequency
dyad and write \(t=\log_D T\). A selector of modulus at most one may depend
on \(\mathfrak k\), the actual equality allocation and the fixed
partner labels. It is independent of the row and the two varying
partner cutoff ideals. In the actual Poisson representation of the
unit part, including the direct term and the corresponding selected
middle subtraction, the target holds in the following regions:
\begin{equation}\label{pca:joint-frequency-frontier}
\begin{array}{c|l}
0\le e\le1/8 & t\le15/32+3e/8\\
1/8\le e\le1/4 & t\le33/64\\
1/4\le e\le5/16 & t\le17/64+e\\
5/16\le e<27/56 & t\le11/32+3e/4\\
e\ge27/56 & \text{all frequency lengths.}
\end{array}
\end{equation}
Thus \(T\lesssim D^{15/32}\) is sufficient over all intermediate
divisors. Constant factors at the endpoints are allowed.
\end{theorem}

For selectors independent of the varying equality labels, further
regions follow by grouping frequencies with prescribed ideal parts.
If the rational content of \(\mathfrak k\) is at least \(D^b\),
with \(0\le b\le9/28\), then \(T\lesssim D^{9/20+3b/5}\)
suffices across all \(d\). Here the content is the largest positive
integer prime to the rational primes under \(S\) whose principal
ideal divides \(\mathfrak k\). At \(E\gtrsim D^{1/4}\),
content of order at least \(D^{41/144}\) permits \(T\lesssim D^{2/3}\).
A good sixth-power base of norm at least of order \(D^{1/16}\),
or an \(S\)-part of norm at least of order \(D^{3/10}\),
suffices at every frequency length. These are conditions on
the actual frequency, with all ideal labels and raw zeros retained.

\begin{theorem}[Equality-wide correction]\label{pca:joint-eqwide}
In the equality-side wide correction, write the equal-norm pair
as \(z_0=g\ell,\ k=g\bar\ell\). Let \(\ell_{\rm g}\) retain the
prime ideals of \(\ell\) whose conjugates also avoid \(S\).
Against the full original partner complement, or its full-frequency
step or wide part with its own subtraction, the target holds on
\(N\ell_{\rm g}\gtrsim D^{\theta(e)}\), where
\[
 \theta(e)=\begin{cases}
 -17/12+37e/6,&17/74\le e\le1/4,\\
 11/24-4e/3,&1/4\le e\le11/32,\\
 0,&\text{otherwise.}
 \end{cases}
\]
In the distinct pure comparable selected rectangles \(B,K\asymp R\)
of the preceding wide analysis, the equality-side prefix extends
to \(Nd\lesssim D^{111/476}\), with the original allocation tests
and the same choices of partner coefficient.
\end{theorem}

The small-conjugate-support band, unbalanced middle divisors and
small-part main frequency range leave contributions to the coherent
signed problem. The full signed assembly and the independent
determinant and zero-location arguments remain open.
